%%%%%%%%%%%%%%%%%%%%%%%%%%%%%%%%%%%%%%%%%%%%%%%%%%%%%%%%%%%%
% 5.13 CALCULUS OF VARIATIONS
% The Composition of Advanced Mathematics
%%%%%%%%%%%%%%%%%%%%%%%%%%%%%%%%%%%%%%%%%%%%%%%%%%%%%%%%%%%%

\section{Calculus of Variations}
\label{sec:calculus-of-variations}

\prerequisite{
Single-variable calculus, multivariable calculus, differential equations,
real analysis, basic functional analysis, integration by parts, and
partial derivatives.
}

\learningobjectives{
By the end of this section, the reader should be able to:
\begin{itemize}
    \item distinguish ordinary optimization from variational optimization;
    \item identify functionals and admissible functions;
    \item construct variations of candidate extremizers;
    \item compute first variations;
    \item derive the Euler--Lagrange equation;
    \item apply the Euler--Lagrange equation to classical variational problems;
    \item handle functionals involving several dependent variables;
    \item recognize higher-order Euler--Lagrange equations;
    \item use natural boundary conditions;
    \item understand constrained variational problems and Lagrange multipliers;
    \item derive the isoperimetric Euler--Lagrange equation;
    \item recognize first integrals such as the Beltrami identity;
    \item compute second variations and interpret stability;
    \item understand convexity and sufficient conditions for minima;
    \item derive Euler--Lagrange PDEs for multiple independent variables;
    \item understand the Dirichlet principle;
    \item recognize variational formulations of differential equations;
    \item understand the basic ideas behind weak formulations and direct methods;
    \item connect symmetry with conservation laws through Noether's theorem;
    \item recognize applications in mechanics, geometry, PDEs, optimization,
          and mathematical physics.
\end{itemize}
}

%%%%%%%%%%%%%%%%%%%%%%%%%%%%%%%%%%%%%%%%%%%%%%%%%%%%%%%%%%%%
\subsection{Optimization Over Functions}
\label{subsec:variational-optimization-functions}
%%%%%%%%%%%%%%%%%%%%%%%%%%%%%%%%%%%%%%%%%%%%%%%%%%%%%%%%%%%%

Ordinary optimization asks for numbers or vectors that minimize or
maximize a function.

For example,

\[
\boxed{
\min_{x\in\R}f(x).
}
\]

The calculus of variations asks for an entire function that optimizes a
quantity.

Schematically,

\[
\boxed{
\min_{y\in\mathcal{A}}J[y].
}
\]

Here \(J\) takes a function as input and produces a number.

The set

\[
\mathcal{A}
\]

is the collection of admissible functions.

Thus calculus of variations is

\[
\boxed{
\text{optimization in spaces of functions}.
}
\]

%%%%%%%%%%%%%%%%%%%%%%%%%%%%%%%%%%%%%%%%%%%%%%%%%%%%%%%%%%%%
\subsection{Functionals}
\label{subsec:functionals}
%%%%%%%%%%%%%%%%%%%%%%%%%%%%%%%%%%%%%%%%%%%%%%%%%%%%%%%%%%%%

\begin{definitionbox}{Functional}
A functional is a mapping that assigns a scalar to a function.

Schematically,

\[
\boxed{
J:\mathcal{A}\to\R.
}
\]

The value of the functional at \(y\) is commonly written

\[
\boxed{
J[y].
}
\]
\end{definitionbox}

A standard variational functional has the form

\[
\boxed{
J[y]
=
\int_a^b
F(x,y(x),y'(x))\,dx.
}
\]

The function \(F\) is often called the \emph{integrand} or
\emph{Lagrangian}.

%%%%%%%%%%%%%%%%%%%%%%%%%%%%%%%%%%%%%%%%%%%%%%%%%%%%%%%%%%%%
\subsection{Functions Versus Functionals}
\label{subsec:functions-versus-functionals}
%%%%%%%%%%%%%%%%%%%%%%%%%%%%%%%%%%%%%%%%%%%%%%%%%%%%%%%%%%%%

A function might map a number to a number:

\[
\boxed{
f:\R\to\R.
}
\]

A functional maps a function to a number:

\[
\boxed{
J:y\mapsto J[y].
}
\]

For example,

\[
J[y]
=
\int_0^1 y(x)^2\,dx
\]

takes the entire function \(y\) as input.

If

\[
y(x)=x,
\]

then

\[
J[y]
=
\int_0^1x^2\,dx
=
\frac13.
\]

%%%%%%%%%%%%%%%%%%%%%%%%%%%%%%%%%%%%%%%%%%%%%%%%%%%%%%%%%%%%
\subsection{Admissible Functions}
\label{subsec:admissible-functions}
%%%%%%%%%%%%%%%%%%%%%%%%%%%%%%%%%%%%%%%%%%%%%%%%%%%%%%%%%%%%

A variational problem does not usually allow arbitrary functions.

The admissible class may impose

\[
\boxed{
\text{boundary conditions},
}
\]

\[
\boxed{
\text{regularity conditions},
}
\]

or

\[
\boxed{
\text{integral constraints}.
}
\]

For example,

\[
\mathcal{A}
=
\{
y\in C^1([a,b]):
y(a)=A,\;
y(b)=B
\}.
\]

The exact choice of function space is an essential part of a rigorous
variational problem.

%%%%%%%%%%%%%%%%%%%%%%%%%%%%%%%%%%%%%%%%%%%%%%%%%%%%%%%%%%%%
\subsection{Classical Variational Problem}
\label{subsec:classical-variational-problem}
%%%%%%%%%%%%%%%%%%%%%%%%%%%%%%%%%%%%%%%%%%%%%%%%%%%%%%%%%%%%

Consider

\[
\boxed{
J[y]
=
\int_a^b
F(x,y,y')\,dx
}
\]

subject to fixed endpoint conditions

\[
\boxed{
y(a)=A,
\qquad
y(b)=B.
}
\]

The goal is to determine a function \(y\) for which \(J[y]\) is minimized
or maximized.

Such a function is called an extremal when it satisfies the necessary
variational equations.

%%%%%%%%%%%%%%%%%%%%%%%%%%%%%%%%%%%%%%%%%%%%%%%%%%%%%%%%%%%%
\subsection{Variations}
\label{subsec:variations}
%%%%%%%%%%%%%%%%%%%%%%%%%%%%%%%%%%%%%%%%%%%%%%%%%%%%%%%%%%%%

Suppose \(y\) is a candidate extremizer.

Perturb \(y\) by defining

\[
\boxed{
y_\varepsilon(x)
=
y(x)+\varepsilon\eta(x).
}
\]

The lowercase Greek letter \(\varepsilon\), read \emph{epsilon}, is a
small real parameter.

The lowercase Greek letter \(\eta\), read \emph{eta}, is the variation
or test function.

For fixed endpoints, require

\[
\boxed{
\eta(a)=\eta(b)=0.
}
\]

Then

\[
y_\varepsilon(a)=A,
\qquad
y_\varepsilon(b)=B
\]

for every sufficiently small \(\varepsilon\).

%%%%%%%%%%%%%%%%%%%%%%%%%%%%%%%%%%%%%%%%%%%%%%%%%%%%%%%%%%%%
\subsection{First Variation}
\label{subsec:first-variation}
%%%%%%%%%%%%%%%%%%%%%%%%%%%%%%%%%%%%%%%%%%%%%%%%%%%%%%%%%%%%

Substitute the perturbed function into the functional:

\[
J[y_\varepsilon]
=
J[y+\varepsilon\eta].
\]

This becomes an ordinary function of the scalar parameter
\(\varepsilon\).

If \(y\) is a local extremizer, then

\[
\boxed{
\left.
\frac{d}{d\varepsilon}
J[y+\varepsilon\eta]
\right|_{\varepsilon=0}
=
0.
}
\]

\begin{definitionbox}{First Variation}
The first variation of \(J\) at \(y\) in the direction \(\eta\) is

\[
\boxed{
\delta J[y;\eta]
=
\left.
\frac{d}{d\varepsilon}
J[y+\varepsilon\eta]
\right|_{\varepsilon=0}.
}
\]
\end{definitionbox}

The lowercase Greek letter \(\delta\), read \emph{delta}, represents a
variation.

%%%%%%%%%%%%%%%%%%%%%%%%%%%%%%%%%%%%%%%%%%%%%%%%%%%%%%%%%%%%
\subsection{Computing the First Variation}
\label{subsec:computing-first-variation}
%%%%%%%%%%%%%%%%%%%%%%%%%%%%%%%%%%%%%%%%%%%%%%%%%%%%%%%%%%%%

Let

\[
J[y]
=
\int_a^b
F(x,y,y')\,dx.
\]

Since

\[
y_\varepsilon
=
y+\varepsilon\eta
\]

and

\[
y_\varepsilon'
=
y'+\varepsilon\eta',
\]

we obtain

\[
J[y_\varepsilon]
=
\int_a^b
F(x,y+\varepsilon\eta,y'+\varepsilon\eta')\,dx.
\]

Differentiating formally with respect to \(\varepsilon\),

\[
\left.
\frac{d}{d\varepsilon}
J[y_\varepsilon]
\right|_{\varepsilon=0}
=
\int_a^b
\left(
F_y\eta
+
F_{y'}\eta'
\right)\,dx.
\]

Thus

\[
\boxed{
\delta J[y;\eta]
=
\int_a^b
\left(
F_y\eta
+
F_{y'}\eta'
\right)\,dx.
}
\]

%%%%%%%%%%%%%%%%%%%%%%%%%%%%%%%%%%%%%%%%%%%%%%%%%%%%%%%%%%%%
\subsection{Integration by Parts}
\label{subsec:variational-integration-by-parts}
%%%%%%%%%%%%%%%%%%%%%%%%%%%%%%%%%%%%%%%%%%%%%%%%%%%%%%%%%%%%

The term involving \(\eta'\) is transformed using integration by parts:

\[
\int_a^b
F_{y'}\eta'\,dx
=
\left[
F_{y'}\eta
\right]_a^b
-
\int_a^b
\frac{d}{dx}
\left(
F_{y'}
\right)\eta\,dx.
\]

Because

\[
\eta(a)=\eta(b)=0,
\]

the boundary term vanishes.

Therefore,

\[
\boxed{
\delta J[y;\eta]
=
\int_a^b
\left[
F_y
-
\frac{d}{dx}
\left(
F_{y'}
\right)
\right]
\eta\,dx.
}
\]

%%%%%%%%%%%%%%%%%%%%%%%%%%%%%%%%%%%%%%%%%%%%%%%%%%%%%%%%%%%%
\subsection{Fundamental Lemma of the Calculus of Variations}
\label{subsec:fundamental-lemma-variations}
%%%%%%%%%%%%%%%%%%%%%%%%%%%%%%%%%%%%%%%%%%%%%%%%%%%%%%%%%%%%

\begin{theorem}[Fundamental Lemma of the Calculus of Variations]
Suppose \(g\) is continuous on \([a,b]\) and

\[
\int_a^b
g(x)\eta(x)\,dx
=
0
\]

for every sufficiently smooth function \(\eta\) satisfying

\[
\eta(a)=\eta(b)=0.
\]

Then

\[
\boxed{
g(x)=0
}
\]

for every \(x\in[a,b]\).
\end{theorem}

The fundamental lemma converts an integral identity valid for every test
function into a pointwise differential equation.

%%%%%%%%%%%%%%%%%%%%%%%%%%%%%%%%%%%%%%%%%%%%%%%%%%%%%%%%%%%%
\subsection{Euler--Lagrange Equation}
\label{subsec:euler-lagrange-equation}
%%%%%%%%%%%%%%%%%%%%%%%%%%%%%%%%%%%%%%%%%%%%%%%%%%%%%%%%%%%%

Applying the fundamental lemma to the first variation gives the central
equation of classical variational calculus.

\begin{theorem}[Euler--Lagrange Equation]
If \(y\) is a sufficiently smooth extremizer of

\[
J[y]
=
\int_a^b
F(x,y,y')\,dx
\]

with fixed endpoints, then

\[
\boxed{
\frac{\partial F}{\partial y}
-
\frac{d}{dx}
\left(
\frac{\partial F}{\partial y'}
\right)
=
0.
}
\]
\end{theorem}

Equivalently,

\[
\boxed{
F_y
-
\frac{d}{dx}F_{y'}
=
0.
}
\]

This differential equation is a necessary condition for an extremum.

%%%%%%%%%%%%%%%%%%%%%%%%%%%%%%%%%%%%%%%%%%%%%%%%%%%%%%%%%%%%
\subsection{Variational Derivative}
\label{subsec:variational-derivative}
%%%%%%%%%%%%%%%%%%%%%%%%%%%%%%%%%%%%%%%%%%%%%%%%%%%%%%%%%%%%

The expression

\[
\boxed{
\frac{\delta J}{\delta y}
=
F_y
-
\frac{d}{dx}F_{y'}
}
\]

is called the variational derivative or functional derivative.

The Euler--Lagrange equation can therefore be written

\[
\boxed{
\frac{\delta J}{\delta y}=0.
}
\]

This resembles the ordinary optimization condition

\[
\boxed{
f'(x)=0.
}
\]

%%%%%%%%%%%%%%%%%%%%%%%%%%%%%%%%%%%%%%%%%%%%%%%%%%%%%%%%%%%%
\subsection{Finite- and Infinite-Dimensional Analogy}
\label{subsec:finite-infinite-variational-analogy}
%%%%%%%%%%%%%%%%%%%%%%%%%%%%%%%%%%%%%%%%%%%%%%%%%%%%%%%%%%%%

Ordinary optimization:

\[
f:\R^n\to\R
\]

has the necessary condition

\[
\boxed{
\nabla f=0.
}
\]

Variational optimization:

\[
J:\mathcal{A}\to\R
\]

has the corresponding condition

\[
\boxed{
\frac{\delta J}{\delta y}=0.
}
\]

Thus the Euler--Lagrange equation acts as an infinite-dimensional version
of the critical-point equation.

%%%%%%%%%%%%%%%%%%%%%%%%%%%%%%%%%%%%%%%%%%%%%%%%%%%%%%%%%%%%
\subsection{Worked Example: Minimum Energy Curve}
\label{subsec:worked-minimum-energy-curve}
%%%%%%%%%%%%%%%%%%%%%%%%%%%%%%%%%%%%%%%%%%%%%%%%%%%%%%%%%%%%

\begin{workedexamplebox}{Minimize the Squared Derivative}

Consider

\[
J[y]
=
\int_0^1
(y')^2\,dx
\]

subject to

\[
y(0)=0,
\qquad
y(1)=1.
\]

The Lagrangian is

\[
F=(y')^2.
\]

Therefore,

\[
F_y=0
\]

and

\[
F_{y'}=2y'.
\]

The Euler--Lagrange equation gives

\[
0-\frac{d}{dx}(2y')=0,
\]

so

\[
y''=0.
\]

Hence

\[
y(x)=C_1x+C_2.
\]

Using

\[
y(0)=0
\]

gives

\[
C_2=0.
\]

Using

\[
y(1)=1
\]

gives

\[
C_1=1.
\]

\begin{answerbox}
\[
\boxed{
y(x)=x.
}
\]
\end{answerbox}

The minimizing curve is the straight line connecting the endpoints.

\end{workedexamplebox}

%%%%%%%%%%%%%%%%%%%%%%%%%%%%%%%%%%%%%%%%%%%%%%%%%%%%%%%%%%%%
\subsection{Shortest Curves in the Plane}
\label{subsec:shortest-plane-curves}
%%%%%%%%%%%%%%%%%%%%%%%%%%%%%%%%%%%%%%%%%%%%%%%%%%%%%%%%%%%%

The length of a graph \(y=y(x)\) is

\[
\boxed{
L[y]
=
\int_a^b
\sqrt{1+(y')^2}\,dx.
}
\]

Thus

\[
F(y')
=
\sqrt{1+(y')^2}.
\]

Since \(F\) does not explicitly depend on \(y\),

\[
F_y=0.
\]

The Euler--Lagrange equation becomes

\[
\frac{d}{dx}
\left(
\frac{y'}{\sqrt{1+(y')^2}}
\right)
=
0.
\]

Hence

\[
\frac{y'}{\sqrt{1+(y')^2}}
=
C,
\]

which implies \(y'\) is constant.

Therefore,

\[
\boxed{
y(x)=mx+b.
}
\]

Straight lines are geodesics of the Euclidean plane.

%%%%%%%%%%%%%%%%%%%%%%%%%%%%%%%%%%%%%%%%%%%%%%%%%%%%%%%%%%%%
\subsection{The Brachistochrone Problem}
\label{subsec:brachistochrone}
%%%%%%%%%%%%%%%%%%%%%%%%%%%%%%%%%%%%%%%%%%%%%%%%%%%%%%%%%%%%

The brachistochrone problem asks:

\[
\boxed{
\text{Which curve gives the shortest travel time under gravity?}
}
\]

The word \emph{brachistochrone} means ``shortest time.''

Suppose a particle begins from rest and descends under constant
gravitational acceleration \(g\).

With an appropriate coordinate convention, conservation of energy gives

\[
v=\sqrt{2gy}.
\]

The travel-time functional becomes

\[
\boxed{
T[y]
=
\int
\frac{\sqrt{1+(y')^2}}
{\sqrt{2gy}}
\,dx.
}
\]

The Euler--Lagrange equation can be simplified using a first integral.

The resulting minimizing curve is an arc of a cycloid.

Thus

\[
\boxed{
\text{fastest path}
\neq
\text{shortest geometric path}.
}
\]

%%%%%%%%%%%%%%%%%%%%%%%%%%%%%%%%%%%%%%%%%%%%%%%%%%%%%%%%%%%%
\subsection{When the Lagrangian Does Not Depend on \(y\)}
\label{subsec:lagrangian-independent-y}
%%%%%%%%%%%%%%%%%%%%%%%%%%%%%%%%%%%%%%%%%%%%%%%%%%%%%%%%%%%%

Suppose

\[
F=F(x,y').
\]

Then

\[
F_y=0.
\]

The Euler--Lagrange equation becomes

\[
\frac{d}{dx}
F_{y'}
=
0.
\]

Therefore,

\[
\boxed{
F_{y'}=C.
}
\]

This provides a first integral of the Euler--Lagrange equation.

%%%%%%%%%%%%%%%%%%%%%%%%%%%%%%%%%%%%%%%%%%%%%%%%%%%%%%%%%%%%
\subsection{Beltrami Identity}
\label{subsec:beltrami-identity}
%%%%%%%%%%%%%%%%%%%%%%%%%%%%%%%%%%%%%%%%%%%%%%%%%%%%%%%%%%%%

Suppose the Lagrangian

\[
F(y,y')
\]

does not explicitly depend on \(x\).

Then the Euler--Lagrange equation implies

\[
\boxed{
F-y'F_{y'}
=
C.
}
\]

This is called the Beltrami identity.

It can reduce a second-order Euler--Lagrange equation to a first-order
equation.

%%%%%%%%%%%%%%%%%%%%%%%%%%%%%%%%%%%%%%%%%%%%%%%%%%%%%%%%%%%%
\subsection{Derivation of the Beltrami Identity}
\label{subsec:beltrami-derivation}
%%%%%%%%%%%%%%%%%%%%%%%%%%%%%%%%%%%%%%%%%%%%%%%%%%%%%%%%%%%%

Differentiate

\[
F-y'F_{y'}.
\]

Since \(F_x=0\),

\[
\frac{dF}{dx}
=
F_yy'
+
F_{y'}y''.
\]

Therefore,

\[
\frac{d}{dx}
\left(
F-y'F_{y'}
\right)
=
F_yy'
+
F_{y'}y''
-
y''F_{y'}
-
y'
\frac{d}{dx}F_{y'}.
\]

Hence

\[
\frac{d}{dx}
\left(
F-y'F_{y'}
\right)
=
y'
\left(
F_y
-
\frac{d}{dx}F_{y'}
\right).
\]

Along an Euler--Lagrange extremal,

\[
F_y-\frac{d}{dx}F_{y'}=0.
\]

Thus

\[
\boxed{
\frac{d}{dx}
\left(
F-y'F_{y'}
\right)
=
0.
}
\]

Therefore,

\[
\boxed{
F-y'F_{y'}=C.
}
\]

%%%%%%%%%%%%%%%%%%%%%%%%%%%%%%%%%%%%%%%%%%%%%%%%%%%%%%%%%%%%
\subsection{Several Dependent Variables}
\label{subsec:several-dependent-variables}
%%%%%%%%%%%%%%%%%%%%%%%%%%%%%%%%%%%%%%%%%%%%%%%%%%%%%%%%%%%%

Suppose

\[
\mathbf{q}(t)
=
(q_1(t),\ldots,q_n(t))
\]

and consider

\[
\boxed{
J[\mathbf{q}]
=
\int_a^b
L(t,\mathbf{q},\dot{\mathbf{q}})\,dt.
}
\]

The dot denotes differentiation with respect to \(t\):

\[
\boxed{
\dot q_j
=
\frac{dq_j}{dt}.
}
\]

The Euler--Lagrange equation holds for each coordinate:

\[
\boxed{
\frac{\partial L}{\partial q_j}
-
\frac{d}{dt}
\left(
\frac{\partial L}{\partial\dot q_j}
\right)
=
0,
\qquad
j=1,\ldots,n.
}
\]

%%%%%%%%%%%%%%%%%%%%%%%%%%%%%%%%%%%%%%%%%%%%%%%%%%%%%%%%%%%%
\subsection{Classical Mechanics}
\label{subsec:variational-classical-mechanics}
%%%%%%%%%%%%%%%%%%%%%%%%%%%%%%%%%%%%%%%%%%%%%%%%%%%%%%%%%%%%

In classical mechanics, define the Lagrangian

\[
\boxed{
L=T-V,
}
\]

where \(T\) is kinetic energy and \(V\) is potential energy.

The action is

\[
\boxed{
S[\mathbf{q}]
=
\int_{t_0}^{t_1}
L(t,\mathbf{q},\dot{\mathbf{q}})\,dt.
}
\]

Hamilton's principle states that physical trajectories are stationary
points of the action:

\[
\boxed{
\delta S=0.
}
\]

The resulting Euler--Lagrange equations are

\[
\boxed{
\frac{d}{dt}
\left(
\frac{\partial L}{\partial\dot q_j}
\right)
-
\frac{\partial L}{\partial q_j}
=
0.
}
\]

%%%%%%%%%%%%%%%%%%%%%%%%%%%%%%%%%%%%%%%%%%%%%%%%%%%%%%%%%%%%
\subsection{Newton's Equation from the Action Principle}
\label{subsec:newton-from-action}
%%%%%%%%%%%%%%%%%%%%%%%%%%%%%%%%%%%%%%%%%%%%%%%%%%%%%%%%%%%%

For a particle in one dimension,

\[
L(x,\dot x)
=
\frac12m\dot x^2-V(x).
\]

Then

\[
\frac{\partial L}{\partial x}
=
-V'(x)
\]

and

\[
\frac{\partial L}{\partial\dot x}
=
m\dot x.
\]

Therefore,

\[
\frac{d}{dt}
\left(
m\dot x
\right)
+
V'(x)
=
0.
\]

Thus

\[
\boxed{
m\ddot x
=
-V'(x).
}
\]

Since force is

\[
F(x)=-V'(x),
\]

we recover

\[
\boxed{
m\ddot x=F.
}
\]

Newton's second law therefore arises from a variational principle.

%%%%%%%%%%%%%%%%%%%%%%%%%%%%%%%%%%%%%%%%%%%%%%%%%%%%%%%%%%%%
\subsection{Conserved Momentum from a Cyclic Coordinate}
\label{subsec:cyclic-coordinate}
%%%%%%%%%%%%%%%%%%%%%%%%%%%%%%%%%%%%%%%%%%%%%%%%%%%%%%%%%%%%

Suppose the Lagrangian does not explicitly depend on a coordinate
\(q_j\).

Then

\[
\frac{\partial L}{\partial q_j}=0.
\]

The Euler--Lagrange equation becomes

\[
\frac{d}{dt}
\left(
\frac{\partial L}{\partial\dot q_j}
\right)
=
0.
\]

Therefore,

\[
\boxed{
\frac{\partial L}{\partial\dot q_j}
=
\text{constant}.
}
\]

The quantity

\[
\boxed{
p_j
=
\frac{\partial L}{\partial\dot q_j}
}
\]

is the generalized momentum conjugate to \(q_j\).

%%%%%%%%%%%%%%%%%%%%%%%%%%%%%%%%%%%%%%%%%%%%%%%%%%%%%%%%%%%%
\subsection{Higher-Order Functionals}
\label{subsec:higher-order-functionals}
%%%%%%%%%%%%%%%%%%%%%%%%%%%%%%%%%%%%%%%%%%%%%%%%%%%%%%%%%%%%

A functional may depend on higher derivatives:

\[
\boxed{
J[y]
=
\int_a^b
F(x,y,y',y'',\ldots,y^{(m)})\,dx.
}
\]

The corresponding Euler--Lagrange equation is

\[
\boxed{
\frac{\partial F}{\partial y}
-
\frac{d}{dx}
\frac{\partial F}{\partial y'}
+
\frac{d^2}{dx^2}
\frac{\partial F}{\partial y''}
-\cdots+
(-1)^m
\frac{d^m}{dx^m}
\frac{\partial F}{\partial y^{(m)}}
=
0.
}
\]

Higher-order variational problems occur in elasticity, beam theory, and
spline interpolation.

%%%%%%%%%%%%%%%%%%%%%%%%%%%%%%%%%%%%%%%%%%%%%%%%%%%%%%%%%%%%
\subsection{Example: Bending Energy}
\label{subsec:bending-energy}
%%%%%%%%%%%%%%%%%%%%%%%%%%%%%%%%%%%%%%%%%%%%%%%%%%%%%%%%%%%%

Consider

\[
\boxed{
J[y]
=
\int_a^b
(y'')^2\,dx.
}
\]

Here

\[
F=(y'')^2.
\]

The higher-order Euler--Lagrange equation gives

\[
\frac{d^2}{dx^2}
\left(
2y''
\right)
=
0.
\]

Therefore,

\[
\boxed{
y^{(4)}=0.
}
\]

Hence the extremals are cubic polynomials.

This is closely related to cubic spline theory and elastic bending.

%%%%%%%%%%%%%%%%%%%%%%%%%%%%%%%%%%%%%%%%%%%%%%%%%%%%%%%%%%%%
\subsection{Variable Endpoints}
\label{subsec:variable-endpoints}
%%%%%%%%%%%%%%%%%%%%%%%%%%%%%%%%%%%%%%%%%%%%%%%%%%%%%%%%%%%%

When endpoint values are not fixed, the variation does not necessarily
satisfy

\[
\eta(a)=\eta(b)=0.
\]

The boundary term

\[
\left[
F_{y'}\eta
\right]_a^b
\]

must then be retained.

This produces additional boundary conditions called natural boundary
conditions or transversality conditions.

%%%%%%%%%%%%%%%%%%%%%%%%%%%%%%%%%%%%%%%%%%%%%%%%%%%%%%%%%%%%
\subsection{Natural Boundary Conditions}
\label{subsec:natural-boundary-conditions}
%%%%%%%%%%%%%%%%%%%%%%%%%%%%%%%%%%%%%%%%%%%%%%%%%%%%%%%%%%%%

Suppose

\[
J[y]
=
\int_a^b
F(x,y,y')\,dx
\]

and \(y(b)\) is free.

The first variation contains the endpoint term

\[
F_{y'}(b)\eta(b).
\]

Since \(\eta(b)\) is arbitrary, stationarity requires

\[
\boxed{
F_{y'}(b)=0.
}
\]

This is a natural boundary condition.

If both endpoint values are free, corresponding conditions arise at both
ends.

%%%%%%%%%%%%%%%%%%%%%%%%%%%%%%%%%%%%%%%%%%%%%%%%%%%%%%%%%%%%
\subsection{Constrained Variational Problems}
\label{subsec:constrained-variational-problems}
%%%%%%%%%%%%%%%%%%%%%%%%%%%%%%%%%%%%%%%%%%%%%%%%%%%%%%%%%%%%

A variational problem may include an integral constraint such as

\[
\boxed{
\int_a^b
G(x,y,y')\,dx
=
C.
}
\]

The problem is then analogous to constrained multivariable optimization.

Introduce a Lagrange multiplier \(\lambda\) and consider

\[
\boxed{
\widetilde{J}[y]
=
\int_a^b
\left(
F+\lambda G
\right)\,dx.
}
\]

The Euler--Lagrange equation is applied to the augmented Lagrangian.

%%%%%%%%%%%%%%%%%%%%%%%%%%%%%%%%%%%%%%%%%%%%%%%%%%%%%%%%%%%%
\subsection{Isoperimetric Problems}
\label{subsec:isoperimetric-problems}
%%%%%%%%%%%%%%%%%%%%%%%%%%%%%%%%%%%%%%%%%%%%%%%%%%%%%%%%%%%%

A classical isoperimetric problem seeks an extremum subject to a fixed
geometric quantity.

For example:

\[
\boxed{
\text{maximize enclosed area subject to fixed perimeter}.
}
\]

The solution is the circle.

In variational form, one quantity becomes the objective functional and
the other becomes an integral constraint.

Such problems motivated much of the historical development of the
calculus of variations.

%%%%%%%%%%%%%%%%%%%%%%%%%%%%%%%%%%%%%%%%%%%%%%%%%%%%%%%%%%%%
\subsection{Second Variation}
\label{subsec:second-variation}
%%%%%%%%%%%%%%%%%%%%%%%%%%%%%%%%%%%%%%%%%%%%%%%%%%%%%%%%%%%%

The first variation provides a necessary stationarity condition.

To distinguish minima from maxima, consider the second variation.

\begin{definitionbox}{Second Variation}
The second variation is

\[
\boxed{
\delta^2J[y;\eta]
=
\left.
\frac{d^2}{d\varepsilon^2}
J[y+\varepsilon\eta]
\right|_{\varepsilon=0}.
}
\]
\end{definitionbox}

For

\[
J[y]
=
\int_a^b
F(x,y,y')\,dx,
\]

the formal second variation is

\[
\boxed{
\delta^2J[y;\eta]
=
\int_a^b
\left(
F_{yy}\eta^2
+
2F_{yy'}\eta\eta'
+
F_{y'y'}(\eta')^2
\right)\,dx.
}
\]

%%%%%%%%%%%%%%%%%%%%%%%%%%%%%%%%%%%%%%%%%%%%%%%%%%%%%%%%%%%%
\subsection{Legendre Necessary Condition}
\label{subsec:legendre-condition}
%%%%%%%%%%%%%%%%%%%%%%%%%%%%%%%%%%%%%%%%%%%%%%%%%%%%%%%%%%%%

For a sufficiently regular weak local minimum, a classical necessary
condition is

\[
\boxed{
F_{y'y'}\geq0
}
\]

along the extremal.

For a strong positivity condition one seeks

\[
\boxed{
F_{y'y'}>0.
}
\]

This is called the strengthened Legendre condition.

It expresses convexity of the Lagrangian with respect to the derivative
variable.

%%%%%%%%%%%%%%%%%%%%%%%%%%%%%%%%%%%%%%%%%%%%%%%%%%%%%%%%%%%%
\subsection{Convexity and Global Minima}
\label{subsec:variational-convexity}
%%%%%%%%%%%%%%%%%%%%%%%%%%%%%%%%%%%%%%%%%%%%%%%%%%%%%%%%%%%%

Convexity can convert stationary points into global minimizers.

Suppose the functional \(J\) is convex on a convex admissible set.

Then any point satisfying the appropriate first-order stationarity
condition is a global minimizer.

Schematically,

\[
\boxed{
\text{convexity}
+
\text{stationarity}
\Longrightarrow
\text{global minimum}.
}
\]

Strict convexity can imply uniqueness.

%%%%%%%%%%%%%%%%%%%%%%%%%%%%%%%%%%%%%%%%%%%%%%%%%%%%%%%%%%%%
\subsection{Example of Strict Convexity}
\label{subsec:strict-convexity-example}
%%%%%%%%%%%%%%%%%%%%%%%%%%%%%%%%%%%%%%%%%%%%%%%%%%%%%%%%%%%%

Consider

\[
J[y]
=
\int_0^1
(y')^2\,dx
\]

with fixed endpoints.

Since

\[
p\mapsto p^2
\]

is strictly convex, the energy is strictly convex in \(y'\).

Therefore the Euler--Lagrange extremal

\[
y(x)=x
\]

is not merely stationary: it is the unique minimizer among the
appropriate admissible functions.

%%%%%%%%%%%%%%%%%%%%%%%%%%%%%%%%%%%%%%%%%%%%%%%%%%%%%%%%%%%%
\subsection{Jacobi Equation Preview}
\label{subsec:jacobi-equation-preview}
%%%%%%%%%%%%%%%%%%%%%%%%%%%%%%%%%%%%%%%%%%%%%%%%%%%%%%%%%%%%

The second variation leads to a differential equation known as the
Jacobi equation.

Its solutions describe infinitesimal variations around an extremal.

Special points called conjugate points can signal loss of minimizing
behavior.

These ideas become especially important in differential geometry, where
they describe variations of geodesics.

%%%%%%%%%%%%%%%%%%%%%%%%%%%%%%%%%%%%%%%%%%%%%%%%%%%%%%%%%%%%
\subsection{Multiple Independent Variables}
\label{subsec:multiple-independent-variables}
%%%%%%%%%%%%%%%%%%%%%%%%%%%%%%%%%%%%%%%%%%%%%%%%%%%%%%%%%%%%

Variational problems may involve functions

\[
u:\Omega\subseteq\R^n\to\R.
\]

Consider

\[
\boxed{
J[u]
=
\int_\Omega
F(x,u,\nabla u)\,dx.
}
\]

The gradient is

\[
\boxed{
\nabla u
=
\left(
\frac{\partial u}{\partial x_1},
\ldots,
\frac{\partial u}{\partial x_n}
\right).
}
\]

The Euler--Lagrange equation becomes a partial differential equation.

%%%%%%%%%%%%%%%%%%%%%%%%%%%%%%%%%%%%%%%%%%%%%%%%%%%%%%%%%%%%
\subsection{Euler--Lagrange PDE}
\label{subsec:euler-lagrange-pde}
%%%%%%%%%%%%%%%%%%%%%%%%%%%%%%%%%%%%%%%%%%%%%%%%%%%%%%%%%%%%

For

\[
J[u]
=
\int_\Omega
F(x,u,\nabla u)\,dx,
\]

the Euler--Lagrange equation is

\[
\boxed{
\frac{\partial F}{\partial u}
-
\sum_{j=1}^{n}
\frac{\partial}{\partial x_j}
\left(
\frac{\partial F}
{\partial u_{x_j}}
\right)
=
0.
}
\]

Using divergence notation,

\[
\boxed{
F_u
-
\nabla\cdot
F_{\nabla u}
=
0.
}
\]

This is one of the fundamental links between variational calculus and
partial differential equations.

%%%%%%%%%%%%%%%%%%%%%%%%%%%%%%%%%%%%%%%%%%%%%%%%%%%%%%%%%%%%
\subsection{Dirichlet Energy}
\label{subsec:dirichlet-energy}
%%%%%%%%%%%%%%%%%%%%%%%%%%%%%%%%%%%%%%%%%%%%%%%%%%%%%%%%%%%%

Consider the Dirichlet energy

\[
\boxed{
E[u]
=
\frac12
\int_\Omega
|\nabla u|^2\,dx.
}
\]

Here

\[
F(\nabla u)
=
\frac12|\nabla u|^2.
\]

Since

\[
F_u=0
\]

and

\[
F_{\nabla u}
=
\nabla u,
\]

the Euler--Lagrange equation becomes

\[
-\nabla\cdot(\nabla u)=0.
\]

Therefore,

\[
\boxed{
\Delta u=0.
}
\]

Thus harmonic functions arise as critical points of Dirichlet energy.

%%%%%%%%%%%%%%%%%%%%%%%%%%%%%%%%%%%%%%%%%%%%%%%%%%%%%%%%%%%%
\subsection{Dirichlet Principle}
\label{subsec:dirichlet-principle}
%%%%%%%%%%%%%%%%%%%%%%%%%%%%%%%%%%%%%%%%%%%%%%%%%%%%%%%%%%%%

The Dirichlet principle states, in an appropriate rigorous setting, that
solutions of certain Laplace boundary-value problems minimize the
Dirichlet energy

\[
\boxed{
E[u]
=
\frac12
\int_\Omega
|\nabla u|^2\,dx
}
\]

among functions satisfying the prescribed boundary values.

Thus

\[
\boxed{
\text{harmonic equation}
\longleftrightarrow
\text{energy minimization}.
}
\]

This principle is central to potential theory and elliptic PDEs.

%%%%%%%%%%%%%%%%%%%%%%%%%%%%%%%%%%%%%%%%%%%%%%%%%%%%%%%%%%%%
\subsection{Poisson Equation from an Energy Functional}
\label{subsec:poisson-variational}
%%%%%%%%%%%%%%%%%%%%%%%%%%%%%%%%%%%%%%%%%%%%%%%%%%%%%%%%%%%%

Consider

\[
\boxed{
J[u]
=
\int_\Omega
\left(
\frac12|\nabla u|^2
-
fu
\right)\,dx.
}
\]

Then

\[
F_u=-f
\]

and

\[
F_{\nabla u}=\nabla u.
\]

The Euler--Lagrange equation gives

\[
-f-\Delta u=0.
\]

Therefore,

\[
\boxed{
-\Delta u=f.
}
\]

The Poisson equation is therefore the Euler--Lagrange equation of an
energy functional.

%%%%%%%%%%%%%%%%%%%%%%%%%%%%%%%%%%%%%%%%%%%%%%%%%%%%%%%%%%%%
\subsection{Minimal Surface Equation}
\label{subsec:minimal-surface-equation}
%%%%%%%%%%%%%%%%%%%%%%%%%%%%%%%%%%%%%%%%%%%%%%%%%%%%%%%%%%%%

For a graph

\[
z=u(x,y),
\]

the surface area functional is

\[
\boxed{
A[u]
=
\int_\Omega
\sqrt{1+|\nabla u|^2}\,dx\,dy.
}
\]

The Euler--Lagrange equation is

\[
\boxed{
\nabla\cdot
\left(
\frac{\nabla u}
{\sqrt{1+|\nabla u|^2}}
\right)
=
0.
}
\]

Solutions describe minimal surfaces.

Thus minimal surface theory is fundamentally variational.

%%%%%%%%%%%%%%%%%%%%%%%%%%%%%%%%%%%%%%%%%%%%%%%%%%%%%%%%%%%%
\subsection{Weak Formulations}
\label{subsec:variational-weak-formulations}
%%%%%%%%%%%%%%%%%%%%%%%%%%%%%%%%%%%%%%%%%%%%%%%%%%%%%%%%%%%%

Suppose

\[
-\Delta u=f
\]

in a domain \(\Omega\).

Multiply by a test function \(\varphi\) and integrate:

\[
\int_\Omega
(-\Delta u)\varphi\,dx
=
\int_\Omega
f\varphi\,dx.
\]

After integration by parts, with suitable boundary conditions,

\[
\boxed{
\int_\Omega
\nabla u\cdot\nabla\varphi\,dx
=
\int_\Omega
f\varphi\,dx.
}
\]

This equation may make sense even when \(u\) does not possess classical
second derivatives.

It is called a weak formulation.

%%%%%%%%%%%%%%%%%%%%%%%%%%%%%%%%%%%%%%%%%%%%%%%%%%%%%%%%%%%%
\subsection{Test Functions}
\label{subsec:variational-test-functions}
%%%%%%%%%%%%%%%%%%%%%%%%%%%%%%%%%%%%%%%%%%%%%%%%%%%%%%%%%%%%

A test function is a sufficiently regular function used to probe an
equation.

A common space is

\[
\boxed{
C_c^\infty(\Omega),
}
\]

the infinitely differentiable functions with compact support contained
in \(\Omega\).

The notation is read

\[
\boxed{
\text{``smooth compactly supported functions on }\Omega\text{.''}
}
\]

Test functions become fundamental in Distribution Theory.

%%%%%%%%%%%%%%%%%%%%%%%%%%%%%%%%%%%%%%%%%%%%%%%%%%%%%%%%%%%%
\subsection{Sobolev-Space Viewpoint}
\label{subsec:variational-sobolev-viewpoint}
%%%%%%%%%%%%%%%%%%%%%%%%%%%%%%%%%%%%%%%%%%%%%%%%%%%%%%%%%%%%

Classical minimizers need not always exist in spaces such as

\[
C^1(\Omega).
\]

Variational problems are therefore often formulated in Sobolev spaces.

For example,

\[
\boxed{
H^1(\Omega)
=
W^{1,2}(\Omega)
}
\]

contains square-integrable functions whose first weak derivatives are
also square integrable.

The Dirichlet energy

\[
\int_\Omega|\nabla u|^2\,dx
\]

is naturally defined on \(H^1(\Omega)\).

%%%%%%%%%%%%%%%%%%%%%%%%%%%%%%%%%%%%%%%%%%%%%%%%%%%%%%%%%%%%
\subsection{Direct Method of the Calculus of Variations}
\label{subsec:direct-method}
%%%%%%%%%%%%%%%%%%%%%%%%%%%%%%%%%%%%%%%%%%%%%%%%%%%%%%%%%%%%

The direct method is a fundamental existence strategy.

To minimize

\[
J[u]
\]

over an admissible set \(\mathcal{A}\):

\[
\boxed{
\text{Step 1: choose a minimizing sequence;}
}
\]

\[
\boxed{
\text{Step 2: obtain compactness or weak compactness;}
}
\]

\[
\boxed{
\text{Step 3: extract a convergent subsequence;}
}
\]

\[
\boxed{
\text{Step 4: show the limit remains admissible;}
}
\]

\[
\boxed{
\text{Step 5: use lower semicontinuity to prove minimality.}
}
\]

This method proves existence without explicitly solving the
Euler--Lagrange equation.

%%%%%%%%%%%%%%%%%%%%%%%%%%%%%%%%%%%%%%%%%%%%%%%%%%%%%%%%%%%%
\subsection{Minimizing Sequences}
\label{subsec:minimizing-sequences}
%%%%%%%%%%%%%%%%%%%%%%%%%%%%%%%%%%%%%%%%%%%%%%%%%%%%%%%%%%%%

Let

\[
m
=
\inf_{u\in\mathcal{A}}J[u].
\]

A sequence

\[
(u_n)\subseteq\mathcal{A}
\]

is a minimizing sequence if

\[
\boxed{
J[u_n]\to m.
}
\]

The central question is whether a subsequence converges to some

\[
u\in\mathcal{A}
\]

and whether

\[
\boxed{
J[u]=m.
}
\]

%%%%%%%%%%%%%%%%%%%%%%%%%%%%%%%%%%%%%%%%%%%%%%%%%%%%%%%%%%%%
\subsection{Coercivity}
\label{subsec:variational-coercivity}
%%%%%%%%%%%%%%%%%%%%%%%%%%%%%%%%%%%%%%%%%%%%%%%%%%%%%%%%%%%%

A functional is coercive when its value becomes large as the norm of the
input becomes large.

Schematically,

\[
\boxed{
\|u\|\to\infty
\quad\Longrightarrow\quad
J[u]\to\infty.
}
\]

Coercivity helps prevent minimizing sequences from escaping to infinity.

It therefore provides boundedness needed for compactness arguments.

%%%%%%%%%%%%%%%%%%%%%%%%%%%%%%%%%%%%%%%%%%%%%%%%%%%%%%%%%%%%
\subsection{Lower Semicontinuity}
\label{subsec:variational-lower-semicontinuity}
%%%%%%%%%%%%%%%%%%%%%%%%%%%%%%%%%%%%%%%%%%%%%%%%%%%%%%%%%%%%

A functional \(J\) is sequentially lower semicontinuous with respect to a
chosen convergence if

\[
u_n\to u
\]

implies

\[
\boxed{
J[u]
\leq
\liminf_{n\to\infty}J[u_n].
}
\]

For weak convergence, one seeks weak lower semicontinuity.

This property ensures that energy does not suddenly decrease when taking
limits of minimizing sequences.

%%%%%%%%%%%%%%%%%%%%%%%%%%%%%%%%%%%%%%%%%%%%%%%%%%%%%%%%%%%%
\subsection{Weak Compactness}
\label{subsec:variational-weak-compactness}
%%%%%%%%%%%%%%%%%%%%%%%%%%%%%%%%%%%%%%%%%%%%%%%%%%%%%%%%%%%%

Infinite-dimensional closed bounded sets are generally not compact in
the norm topology.

However, reflexive Banach spaces often provide weak compactness for
bounded sequences.

Thus one frequently obtains

\[
\boxed{
u_{n_k}
\rightharpoonup
u,
}
\]

where the symbol

\[
\rightharpoonup
\]

denotes weak convergence.

Weak convergence is one of the central tools of modern variational
analysis.

%%%%%%%%%%%%%%%%%%%%%%%%%%%%%%%%%%%%%%%%%%%%%%%%%%%%%%%%%%%%
\subsection{Existence by the Direct Method}
\label{subsec:existence-direct-method}
%%%%%%%%%%%%%%%%%%%%%%%%%%%%%%%%%%%%%%%%%%%%%%%%%%%%%%%%%%%%

A common theorem pattern is:

\begin{theorem}[Direct Method Principle]
Suppose the admissible set is nonempty and weakly closed in a reflexive
Banach space. Suppose further that the functional \(J\) is coercive and
weakly lower semicontinuous.

Then, under the corresponding standard hypotheses, \(J\) attains its
minimum on the admissible set.
\end{theorem}

This principle underlies many existence results in PDEs and optimization.

%%%%%%%%%%%%%%%%%%%%%%%%%%%%%%%%%%%%%%%%%%%%%%%%%%%%%%%%%%%%
\subsection{Noether's Theorem}
\label{subsec:noethers-theorem}
%%%%%%%%%%%%%%%%%%%%%%%%%%%%%%%%%%%%%%%%%%%%%%%%%%%%%%%%%%%%

A deep connection exists between symmetries of a variational problem and
conservation laws.

\begin{theorem}[Noether's Theorem --- Informal Form]
Every sufficiently smooth continuous symmetry of an action functional
produces a corresponding conserved quantity along Euler--Lagrange
trajectories.
\end{theorem}

Schematically,

\[
\boxed{
\text{continuous symmetry}
\Longrightarrow
\text{conservation law}.
}
\]

%%%%%%%%%%%%%%%%%%%%%%%%%%%%%%%%%%%%%%%%%%%%%%%%%%%%%%%%%%%%
\subsection{Examples of Noether Symmetries}
\label{subsec:noether-examples}
%%%%%%%%%%%%%%%%%%%%%%%%%%%%%%%%%%%%%%%%%%%%%%%%%%%%%%%%%%%%

In classical mechanics:

\[
\boxed{
\text{time-translation symmetry}
\Longrightarrow
\text{energy conservation},
}
\]

\[
\boxed{
\text{spatial-translation symmetry}
\Longrightarrow
\text{linear momentum conservation},
}
\]

and

\[
\boxed{
\text{rotational symmetry}
\Longrightarrow
\text{angular momentum conservation}.
}
\]

Noether's theorem explains these laws as manifestations of symmetry
rather than unrelated physical coincidences.

%%%%%%%%%%%%%%%%%%%%%%%%%%%%%%%%%%%%%%%%%%%%%%%%%%%%%%%%%%%%
\subsection{Hamiltonian Connection}
\label{subsec:variational-hamiltonian}
%%%%%%%%%%%%%%%%%%%%%%%%%%%%%%%%%%%%%%%%%%%%%%%%%%%%%%%%%%%%

For a Lagrangian

\[
L(q,\dot q,t),
\]

define generalized momentum

\[
\boxed{
p
=
\frac{\partial L}{\partial\dot q}.
}
\]

When the Legendre transformation is nondegenerate, define the Hamiltonian

\[
\boxed{
H(q,p,t)
=
p\dot q-L.
}
\]

The Euler--Lagrange equations can then be transformed into Hamilton's
equations

\[
\boxed{
\dot q
=
\frac{\partial H}{\partial p},
}
\]

\[
\boxed{
\dot p
=
-\frac{\partial H}{\partial q}.
}
\]

This connects variational calculus with symplectic geometry and
Hamiltonian mechanics.

%%%%%%%%%%%%%%%%%%%%%%%%%%%%%%%%%%%%%%%%%%%%%%%%%%%%%%%%%%%%
\subsection{Geodesics as Variational Problems}
\label{subsec:geodesics-variational}
%%%%%%%%%%%%%%%%%%%%%%%%%%%%%%%%%%%%%%%%%%%%%%%%%%%%%%%%%%%%

On a Riemannian manifold, a curve

\[
\gamma:[a,b]\to M
\]

has length

\[
\boxed{
L[\gamma]
=
\int_a^b
\|\dot\gamma(t)\|\,dt.
}
\]

One may instead study the energy

\[
\boxed{
E[\gamma]
=
\frac12
\int_a^b
\|\dot\gamma(t)\|^2\,dt.
}
\]

Critical points satisfy the geodesic equation.

Thus geodesics generalize the Euclidean fact that straight lines minimize
distance locally.

%%%%%%%%%%%%%%%%%%%%%%%%%%%%%%%%%%%%%%%%%%%%%%%%%%%%%%%%%%%%
\subsection{Eigenvalues from Variational Principles}
\label{subsec:eigenvalues-variational}
%%%%%%%%%%%%%%%%%%%%%%%%%%%%%%%%%%%%%%%%%%%%%%%%%%%%%%%%%%%%

Eigenvalue problems can often be reformulated variationally.

For example, under suitable boundary conditions, the smallest eigenvalue
of the Laplacian can be characterized using the Rayleigh quotient

\[
\boxed{
R[u]
=
\frac{
\int_\Omega|\nabla u|^2\,dx
}{
\int_\Omega|u|^2\,dx
}.
}
\]

One seeks

\[
\boxed{
\lambda_1
=
\inf_{u\neq0}R[u].
}
\]

This connects calculus of variations directly to operator theory and
spectral theory.

%%%%%%%%%%%%%%%%%%%%%%%%%%%%%%%%%%%%%%%%%%%%%%%%%%%%%%%%%%%%
\subsection{Rayleigh Quotient}
\label{subsec:rayleigh-quotient}
%%%%%%%%%%%%%%%%%%%%%%%%%%%%%%%%%%%%%%%%%%%%%%%%%%%%%%%%%%%%

For a self-adjoint matrix or operator \(A\), the Rayleigh quotient is

\[
\boxed{
R[u]
=
\frac{\langle Au,u\rangle}
{\langle u,u\rangle},
\qquad
u\neq0.
}
\]

In finite dimensions, minimizing and maximizing this quotient produces
extreme eigenvalues of a Hermitian matrix.

The infinite-dimensional version provides variational characterizations
of eigenvalues of many differential operators.

%%%%%%%%%%%%%%%%%%%%%%%%%%%%%%%%%%%%%%%%%%%%%%%%%%%%%%%%%%%%
\subsection{Gradient Flows}
\label{subsec:gradient-flows}
%%%%%%%%%%%%%%%%%%%%%%%%%%%%%%%%%%%%%%%%%%%%%%%%%%%%%%%%%%%%

Instead of solving

\[
\frac{\delta J}{\delta u}=0
\]

directly, one may evolve a function in the direction of decreasing
energy:

\[
\boxed{
\frac{\partial u}{\partial t}
=
-
\frac{\delta J}{\delta u}.
}
\]

This is called a gradient flow.

Ideally,

\[
\boxed{
\frac{d}{dt}J[u(t)]
\leq0.
}
\]

Gradient flows connect optimization with time-dependent differential
equations.

%%%%%%%%%%%%%%%%%%%%%%%%%%%%%%%%%%%%%%%%%%%%%%%%%%%%%%%%%%%%
\subsection{Heat Equation as Gradient Flow}
\label{subsec:heat-gradient-flow}
%%%%%%%%%%%%%%%%%%%%%%%%%%%%%%%%%%%%%%%%%%%%%%%%%%%%%%%%%%%%

For the Dirichlet energy

\[
E[u]
=
\frac12
\int_\Omega
|\nabla u|^2\,dx,
\]

the variational derivative is formally

\[
\boxed{
\frac{\delta E}{\delta u}
=
-\Delta u.
}
\]

Therefore the \(L^2\)-gradient flow is

\[
\frac{\partial u}{\partial t}
=
-\frac{\delta E}{\delta u}.
\]

Hence

\[
\boxed{
u_t=\Delta u.
}
\]

Thus the heat equation can be interpreted as an energy-dissipation
process.

%%%%%%%%%%%%%%%%%%%%%%%%%%%%%%%%%%%%%%%%%%%%%%%%%%%%%%%%%%%%
\subsection{Variational Inequalities Preview}
\label{subsec:variational-inequalities-preview}
%%%%%%%%%%%%%%%%%%%%%%%%%%%%%%%%%%%%%%%%%%%%%%%%%%%%%%%%%%%%

If the admissible set contains inequality constraints, the Euler--Lagrange
equation may be replaced by a variational inequality.

A typical form is

\[
\boxed{
a(u,v-u)
\geq
\ell(v-u)
}
\]

for every admissible \(v\).

Variational inequalities arise in

\[
\boxed{
\text{obstacle problems},
\quad
\text{contact mechanics},
\quad
\text{constrained optimization}.
}
\]

%%%%%%%%%%%%%%%%%%%%%%%%%%%%%%%%%%%%%%%%%%%%%%%%%%%%%%%%%%%%
\subsection{Optimal Control Connection}
\label{subsec:variational-optimal-control}
%%%%%%%%%%%%%%%%%%%%%%%%%%%%%%%%%%%%%%%%%%%%%%%%%%%%%%%%%%%%

Optimal control generalizes variational calculus by allowing a system to
evolve according to a differential equation

\[
\boxed{
\dot x=f(x,u,t),
}
\]

where \(u\) is a control.

One seeks a control minimizing a cost such as

\[
\boxed{
J[u]
=
\int_{t_0}^{t_1}
L(x(t),u(t),t)\,dt.
}
\]

The calculus of variations provides the conceptual foundation for
Pontryagin's Maximum Principle and modern optimal control.

%%%%%%%%%%%%%%%%%%%%%%%%%%%%%%%%%%%%%%%%%%%%%%%%%%%%%%%%%%%%
\subsection{Common Errors}
\label{subsec:calculus-variations-common-errors}
%%%%%%%%%%%%%%%%%%%%%%%%%%%%%%%%%%%%%%%%%%%%%%%%%%%%%%%%%%%%

\subsubsection{Confusing a Function with a Functional}

A function maps points to values.

A functional maps functions to values.

%%%%%%%%%%%%%%%%%%%%%%%%%%%%%%%%%%%%%%%%%%%%%%%%%%%%%%%%%%%%
\subsubsection{Forgetting Endpoint Conditions on Variations}

For fixed endpoints,

\[
\boxed{
\eta(a)=\eta(b)=0.
}
\]

Without this condition, the boundary terms from integration by parts do
not vanish.

%%%%%%%%%%%%%%%%%%%%%%%%%%%%%%%%%%%%%%%%%%%%%%%%%%%%%%%%%%%%
\subsubsection{Forgetting Integration by Parts}

The first variation initially contains

\[
F_{y'}\eta'.
\]

The fundamental lemma cannot be applied directly until derivatives are
transferred away from the arbitrary variation.

%%%%%%%%%%%%%%%%%%%%%%%%%%%%%%%%%%%%%%%%%%%%%%%%%%%%%%%%%%%%
\subsubsection{Treating the Euler--Lagrange Equation as Sufficient}

The Euler--Lagrange equation gives a necessary condition under appropriate
hypotheses.

A solution may represent

\[
\boxed{
\text{a minimum},
\quad
\text{a maximum},
\quad
\text{or another stationary point}.
}
\]

Additional analysis is required.

%%%%%%%%%%%%%%%%%%%%%%%%%%%%%%%%%%%%%%%%%%%%%%%%%%%%%%%%%%%%
\subsubsection{Ignoring the Admissible Function Space}

A formal Euler--Lagrange solution may fail to satisfy

\[
\boxed{
\text{boundary conditions},
\quad
\text{regularity},
\quad
\text{constraints}.
}
\]

Such a function is not an admissible solution.

%%%%%%%%%%%%%%%%%%%%%%%%%%%%%%%%%%%%%%%%%%%%%%%%%%%%%%%%%%%%
\subsubsection{Assuming a Minimizer Exists}

An infimum need not be attained.

Existence often requires compactness, coercivity, lower semicontinuity,
or other structural hypotheses.

%%%%%%%%%%%%%%%%%%%%%%%%%%%%%%%%%%%%%%%%%%%%%%%%%%%%%%%%%%%%
\subsubsection{Confusing Weak and Classical Solutions}

A weak solution may not possess enough derivatives to satisfy the
differential equation pointwise.

It satisfies an integrated test-function formulation instead.

%%%%%%%%%%%%%%%%%%%%%%%%%%%%%%%%%%%%%%%%%%%%%%%%%%%%%%%%%%%%
\subsubsection{Ignoring Boundary Terms}

Boundary terms may encode physically and mathematically meaningful
natural boundary conditions.

They should not be discarded unless the variations actually vanish there.

%%%%%%%%%%%%%%%%%%%%%%%%%%%%%%%%%%%%%%%%%%%%%%%%%%%%%%%%%%%%
\subsubsection{Using the Beltrami Identity When \(F\) Depends on \(x\)}

The identity

\[
F-y'F_{y'}=C
\]

requires that \(F\) have no explicit dependence on \(x\).

%%%%%%%%%%%%%%%%%%%%%%%%%%%%%%%%%%%%%%%%%%%%%%%%%%%%%%%%%%%%
\subsubsection{Confusing Stationarity with Stability}

The first variation detects stationarity.

The second variation and related conditions provide information about
local minimizing or maximizing behavior.

%%%%%%%%%%%%%%%%%%%%%%%%%%%%%%%%%%%%%%%%%%%%%%%%%%%%%%%%%%%%
\subsubsection{Ignoring Function-Space Topology}

In modern variational problems, whether a functional is continuous,
compact, or lower semicontinuous depends on the topology being used.

Strong and weak convergence can produce very different conclusions.

%%%%%%%%%%%%%%%%%%%%%%%%%%%%%%%%%%%%%%%%%%%%%%%%%%%%%%%%%%%%
\subsection{Applications}
\label{subsec:calculus-variations-applications}
%%%%%%%%%%%%%%%%%%%%%%%%%%%%%%%%%%%%%%%%%%%%%%%%%%%%%%%%%%%%

\begin{applicationbox}
\textbf{Classical Mechanics.}
Hamilton's principle derives equations of motion from stationary action.

\medskip

\textbf{Geometry.}
Geodesics arise as critical points of length or energy functionals.

\medskip

\textbf{Minimal Surfaces.}
Surface-area minimization leads to nonlinear Euler--Lagrange PDEs.

\medskip

\textbf{Partial Differential Equations.}
Many elliptic equations are Euler--Lagrange equations of energy
functionals.

\medskip

\textbf{Potential Theory.}
Harmonic functions minimize Dirichlet energy under fixed boundary data.

\medskip

\textbf{Elasticity.}
Equilibrium configurations minimize elastic energy.

\medskip

\textbf{Fluid Mechanics.}
Variational and energy methods provide formulations and estimates for
fluid equations.

\medskip

\textbf{General Relativity.}
Einstein's field equations arise from variation of the Einstein--Hilbert
action.

\medskip

\textbf{Quantum Field Theory.}
Field equations are derived from action functionals involving fields and
their derivatives.

\medskip

\textbf{Optimal Control.}
Variational ideas lead to necessary conditions for optimal trajectories
and controls.

\medskip

\textbf{Image Processing.}
Variational models minimize energies balancing image fidelity and
regularity.

\medskip

\textbf{Machine Learning.}
Optimization over functions, regularization, and gradient-flow viewpoints
have close connections with variational principles.

\medskip

\textbf{Spectral Theory.}
Rayleigh quotients and minimax principles characterize eigenvalues of
self-adjoint operators.

\medskip

\textbf{Numerical Analysis.}
Finite-element methods often begin from variational or weak formulations
of differential equations.
\end{applicationbox}

%%%%%%%%%%%%%%%%%%%%%%%%%%%%%%%%%%%%%%%%%%%%%%%%%%%%%%%%%%%%
\subsection{Exercises}
\label{subsec:calculus-variations-exercises}
%%%%%%%%%%%%%%%%%%%%%%%%%%%%%%%%%%%%%%%%%%%%%%%%%%%%%%%%%%%%

\begin{exercise}[1]
Determine whether each object is a function or functional:

\[
f(x)=x^2,
\]

\[
J[y]=\int_0^1y(x)^2\,dx.
\]
\end{exercise}

\begin{exercise}[1]
For

\[
J[y]
=
\int_0^1
(y')^2\,dx,
\]

compute the first variation.
\end{exercise}

\begin{exercise}[1]
Explain why fixed endpoint conditions imply

\[
\eta(a)=\eta(b)=0.
\]
\end{exercise}

\begin{exercise}[1]
Find the Euler--Lagrange equation for

\[
J[y]
=
\int_a^b
\left(
(y')^2+y^2
\right)\,dx.
\]
\end{exercise}

\begin{exercise}[1]
Find the Euler--Lagrange equation for

\[
J[y]
=
\int_a^b
\left(
\frac12(y')^2-V(y)
\right)\,dx.
\]
\end{exercise}

\begin{exercise}[2]
Derive the Euler--Lagrange equation directly from the first variation.
\end{exercise}

\begin{exercise}[2]
Show that the shortest graph joining two fixed points in the plane is a
straight line.
\end{exercise}

\begin{exercise}[2]
Use the Beltrami identity to simplify the Euler--Lagrange equation when

\[
F=F(y,y').
\]
\end{exercise}

\begin{exercise}[2]
Find the extremals of

\[
J[y]
=
\int_0^1
\left(
(y')^2+2y
\right)\,dx
\]

with

\[
y(0)=y(1)=0.
\]
\end{exercise}

\begin{exercise}[2]
For

\[
L(x,\dot x)
=
\frac12m\dot x^2-V(x),
\]

derive Newton's equation.
\end{exercise}

\begin{exercise}[2]
Suppose \(L\) does not depend on \(q_j\). Show that the corresponding
generalized momentum is conserved.
\end{exercise}

\begin{exercise}[2]
Find the higher-order Euler--Lagrange equation for

\[
J[y]
=
\int_a^b
(y'')^2\,dx.
\]
\end{exercise}

\begin{exercise}[2]
Determine the natural boundary condition at \(x=b\) when \(y(b)\) is
free.
\end{exercise}

\begin{exercise}[2]
Compute the second variation of

\[
J[y]
=
\int_a^b
(y')^2\,dx.
\]
\end{exercise}

\begin{exercise}[2]
Show that this second variation is nonnegative.
\end{exercise}

\begin{exercise}[2]
Derive the Euler--Lagrange PDE for

\[
J[u]
=
\int_\Omega
\left(
\frac12|\nabla u|^2-fu
\right)\,dx.
\]
\end{exercise}

\begin{exercise}[2]
Show that the resulting equation is

\[
-\Delta u=f.
\]
\end{exercise}

\begin{exercise}[3]
Derive the minimal surface equation from

\[
A[u]
=
\int_\Omega
\sqrt{1+|\nabla u|^2}\,dx.
\]
\end{exercise}

\begin{exercise}[3]
Derive the weak formulation of the Poisson equation with homogeneous
Dirichlet boundary conditions.
\end{exercise}

\begin{exercise}[3]
Explain why the Dirichlet energy is naturally defined on \(H^1(\Omega)\).
\end{exercise}

\begin{exercise}[3]
Prove that a strictly convex functional has at most one minimizer on a
convex admissible set.
\end{exercise}

\begin{exercise}[3]
Use a Lagrange multiplier to derive the Euler--Lagrange equation for a
functional subject to an integral constraint.
\end{exercise}

\begin{exercise}[3]
Formulate the classical isoperimetric problem as a constrained
variational problem.
\end{exercise}

\begin{exercise}[3]
Derive the brachistochrone functional using conservation of energy.
\end{exercise}

\begin{exercise}[3]
Use the Beltrami identity to reduce the brachistochrone equation to a
first-order equation.
\end{exercise}

\begin{exercise}[3]
Show formally that the \(L^2\)-gradient flow of Dirichlet energy is the
heat equation.
\end{exercise}

\begin{exercise}[3]
Derive the Rayleigh quotient for the first Dirichlet eigenvalue of the
Laplacian.
\end{exercise}

\begin{exercise}[3]
Show that stationary points of the Rayleigh quotient satisfy an
eigenvalue equation under appropriate hypotheses.
\end{exercise}

\begin{exercise}[4]
Prove the fundamental lemma of the calculus of variations.
\end{exercise}

\begin{exercise}[4]
Study sufficient conditions for a weak local minimum using the second
variation.
\end{exercise}

\begin{exercise}[4]
Derive the Jacobi equation associated with the second variation.
\end{exercise}

\begin{exercise}[4]
Study the relationship between conjugate points and minimizing
geodesics.
\end{exercise}

\begin{exercise}[4]
Prove an existence theorem using the direct method in a reflexive Banach
space.
\end{exercise}

\begin{exercise}[4]
Study weak lower semicontinuity of convex integral functionals.
\end{exercise}

\begin{exercise}[4]
Derive a conservation law from time-translation invariance of a
Lagrangian.
\end{exercise}

\begin{exercise}[4]
Study Noether's theorem for a one-parameter family of transformations.
\end{exercise}

\begin{exercise}[4]
Derive Hamilton's equations from the Euler--Lagrange formulation through
the Legendre transform.
\end{exercise}

\begin{exercise}[4]
Study the variational characterization of eigenvalues through the
minimax principle.
\end{exercise}

%%%%%%%%%%%%%%%%%%%%%%%%%%%%%%%%%%%%%%%%%%%%%%%%%%%%%%%%%%%%
\subsection{Conceptual Summary}
\label{subsec:calculus-variations-summary}
%%%%%%%%%%%%%%%%%%%%%%%%%%%%%%%%%%%%%%%%%%%%%%%%%%%%%%%%%%%%

The calculus of variations studies optimization problems whose unknowns
are functions.

A typical functional is

\[
\boxed{
J[y]
=
\int_a^b
F(x,y,y')\,dx.
}
\]

A variation is

\[
\boxed{
y_\varepsilon
=
y+\varepsilon\eta.
}
\]

For fixed endpoints,

\[
\boxed{
\eta(a)=\eta(b)=0.
}
\]

The first variation is

\[
\boxed{
\delta J[y;\eta]
=
\left.
\frac{d}{d\varepsilon}
J[y+\varepsilon\eta]
\right|_{\varepsilon=0}.
}
\]

Stationarity leads to the Euler--Lagrange equation

\[
\boxed{
F_y
-
\frac{d}{dx}F_{y'}
=
0.
}
\]

If \(F\) does not explicitly depend on \(x\), the Beltrami identity gives

\[
\boxed{
F-y'F_{y'}
=
C.
}
\]

For several dependent variables,

\[
\boxed{
\frac{d}{dt}
\left(
\frac{\partial L}{\partial\dot q_j}
\right)
-
\frac{\partial L}{\partial q_j}
=
0.
}
\]

For functionals over \(\R^n\),

\[
J[u]
=
\int_\Omega
F(x,u,\nabla u)\,dx,
\]

the Euler--Lagrange equation becomes

\[
\boxed{
F_u
-
\nabla\cdot F_{\nabla u}
=
0.
}
\]

The Dirichlet energy

\[
\boxed{
E[u]
=
\frac12
\int_\Omega|\nabla u|^2\,dx
}
\]

produces

\[
\boxed{
\Delta u=0.
}
\]

Modern variational analysis uses

\[
\boxed{
\text{weak convergence},
\quad
\text{coercivity},
\quad
\text{lower semicontinuity},
\quad
\text{compactness}
}
\]

to prove that minimizers exist.

Noether's theorem connects

\[
\boxed{
\text{symmetry}
\Longrightarrow
\text{conservation law}.
}
\]

The calculus of variations therefore provides a unified framework for

\[
\boxed{
\text{optimization},
\quad
\text{geometry},
\quad
\text{mechanics},
\quad
\text{PDEs},
\quad
\text{spectral theory}.
}
\]

%%%%%%%%%%%%%%%%%%%%%%%%%%%%%%%%%%%%%%%%%%%%%%%%%%%%%%%%%%%%
\subsection{Section Connections}
\label{subsec:calculus-variations-connections}
%%%%%%%%%%%%%%%%%%%%%%%%%%%%%%%%%%%%%%%%%%%%%%%%%%%%%%%%%%%%

\chapterconnection{
Calculus of Variations extends the optimization ideas of ordinary and
multivariable calculus to infinite-dimensional spaces of functions.
Functional Analysis supplies the Banach-space, Hilbert-space, weak
convergence, and compactness framework required for modern existence
theory. Operator Theory connects variational principles to eigenvalue
problems through Rayleigh quotients and minimax principles. Harmonic
Analysis and Potential Theory use energy minimization to study harmonic
functions and elliptic equations. Differential Geometry interprets
geodesics and minimal surfaces as variational objects. Differential
Equations and PDEs arise as Euler--Lagrange equations and gradient flows.
Symplectic Geometry and Hamiltonian Mechanics are connected through the
Legendre transform, while Noether's theorem links variational symmetry
with conservation laws. These ideas also lead naturally into Distribution
Theory, Geometric Analysis, Optimization, Optimal Control, and
Mathematical Physics.
}

%%%%%%%%%%%%%%%%%%%%%%%%%%%%%%%%%%%%%%%%%%%%%%%%%%%%%%%%%%%%
% SECTION SOURCE TRACKING
%%%%%%%%%%%%%%%%%%%%%%%%%%%%%%%%%%%%%%%%%%%%%%%%%%%%%%%%%%%%

%------------------------------------------------------------
% 5.13 Calculus of Variations
%------------------------------------------------------------
%
% Source basis:
%   - Standard undergraduate and beginning graduate calculus of
%     variations.
%   - Standard variational methods in analysis and PDEs.
%   - Standard Lagrangian mechanics and variational principles.
%
% Used for:
%   - functionals;
%   - admissible functions;
%   - variations;
%   - first variations;
%   - fundamental lemma;
%   - Euler--Lagrange equations;
%   - variational derivatives;
%   - classical extremal problems;
%   - shortest curves;
%   - brachistochrone problem;
%   - first integrals;
%   - Beltrami identity;
%   - several dependent variables;
%   - Lagrangian mechanics;
%   - generalized momentum;
%   - higher-order variational problems;
%   - natural boundary conditions;
%   - constrained variations;
%   - isoperimetric problems;
%   - second variation;
%   - Legendre condition;
%   - convexity;
%   - Jacobi equation preview;
%   - multiple independent variables;
%   - Euler--Lagrange PDEs;
%   - Dirichlet energy;
%   - Dirichlet principle;
%   - Poisson equation;
%   - minimal surfaces;
%   - weak formulations;
%   - test functions;
%   - Sobolev-space viewpoint;
%   - direct method;
%   - minimizing sequences;
%   - coercivity;
%   - lower semicontinuity;
%   - weak compactness;
%   - Noether's theorem;
%   - Hamiltonian connection;
%   - geodesics;
%   - Rayleigh quotients;
%   - gradient flows;
%   - variational inequalities;
%   - optimal control preview.
%
% Notes:
%   - Functional Analysis supplies the infinite-dimensional setting
%     needed for rigorous variational existence theory.
%   - Operator Theory supplies the spectral interpretation of
%     Rayleigh quotients and energy minimization.
%   - Distribution Theory develops weak derivatives and generalized
%     solutions further.
%   - Potential Theory develops the Dirichlet principle and harmonic
%     energy methods in greater depth.
%   - Geometric Analysis develops variational problems involving
%     geodesics, curvature, and minimal surfaces.
%   - Optimization later develops constrained and infinite-dimensional
%     optimization from a broader perspective.
%
%%%%%%%%%%%%%%%%%%%%%%%%%%%%%%%%%%%%%%%%%%%%%%%%%%%%%%%%%%%%